% --------------------------------------------------------------
% SURAJ A — Data Engineer Resume
% Overleaf: New Project > Upload > select this file
% Compiler: pdfLaTeX
% --------------------------------------------------------------
\documentclass[10pt]{article}

\usepackage{geometry}
\geometry{a4paper, margin=0.7in}

\usepackage[T1]{fontenc}
\usepackage{lmodern}
\renewcommand{\familydefault}{\sfdefault}

\usepackage{parskip}
\setlength{\parskip}{3pt}

\usepackage{hyperref}
\hypersetup{colorlinks=true, linkcolor=blue!60!black, urlcolor=blue!60!black}

\usepackage{titlesec}
\titleformat{\section}{\large\bfseries\uppercase}{}{0em}{}[\vspace{-4pt}\rule{\textwidth}{0.5pt}]
\titlespacing{\section}{0pt}{8pt}{4pt}

\titleformat{\subsection}[runin]{\bfseries}{}{0em}{}
\titlespacing{\subsection}{0pt}{4pt}{6pt}

\pagenumbering{gcf}

\begin{document}

% === HEADER ===
\begin{center}
  {\LARGE\textbf{SURAJ A}}\\[4pt]
  \small
  \href{mailto:surajra0628@gmail.com}{surajra0628@gmail.com} $|$
  \href{tel:+919363064502}{+91 9363064502} $|$
  \href{https://linkedin.com/in/suraj314}{linkedin.com/in/suraj314} $|$
  \href{https://github.com/suraj-0628}{github.com/suraj-0628} $|$
  Coimbatore, Tamil Nadu
\end{center}

% === PROFESSIONAL SUMMARY ===
\section{Professional Summary}
Data Engineer with strong foundations in Python, SQL, ETL pipeline design, and database systems. Experienced in building automated data pipelines that scrape, clean, transform, and load data from multiple heterogeneous sources. Skilled in data modeling, query optimization, geocoding infrastructure, and real-time streaming with Kafka. Built production-grade scraper and ETL systems processing over 2,000 records per week with built-in quality checks and monitoring. Eager to design and maintain robust data infrastructure at scale.

% === TECHNICAL SKILLS ===
\section{Technical Skills}
\textbf{Languages:} Python, SQL, Java, Shell Scripting \\
\textbf{Data Pipeline \& ETL:} Apache Airflow, Apache Kafka, APScheduler, httpx, BeautifulSoup, Celery \\
\textbf{Databases \& Storage:} PostgreSQL, SQLite, Redis, Parquet, H2 \\
\textbf{Data Processing:} Pandas, NumPy, Apache Spark, SQLAlchemy, Alembic \\
\textbf{Cloud \& DevOps:} Docker, Google Cloud Compute Engine, Nginx, systemd, Linux \\
\textbf{Monitoring \& Quality:} Data validation (Pydantic), logging (structlog), health checks, anomaly detection \\
\textbf{Concepts:} Star Schema, Lakehouse Architecture (Bronze/Silver/Gold), Change Data Capture, Idempotency, Backfill

% === PROJECTS ===
\section{Projects}

\subsection{Scalable ETL Pipeline for Multi-Source Event Analytics}
\textit{Python, Apache Airflow, PostgreSQL, Redis, Docker, httpx, APScheduler}
\begin{itemize}
  \item Designed and deployed an automated ETL pipeline that scrapes 2,000 college events per week from 10 heterogeneous sources including Knowafest and direct submissions. Uses intelligent deduplication with normalized URL hashing.
  \item Built a multi-tier geocoding microservice backed by Nominatim API with an LRU cache, rate limiting at 1 request per second, and a city-coord fallback dictionary. Resolves 95\% of venue coordinates.
  \item Implemented data quality validation with Pydantic schemas, null checks, date sanity rules, duplicate detection, and automatic Slack alerts on pipeline failures.
  \item Optimized PostgreSQL queries with composite indexes. Reduced event listing latency from 1.2 seconds to 40 milliseconds under 50 concurrent requests.
  \item Containerized with Docker and deployed on Google Cloud Platform with a systemd managed service supporting zero-downtime redeploys.
  \item Used incremental processing with watermark tracking. Avoids full re-scrapes and only processes new or changed events.
\end{itemize}

\subsection{Real-Time Data Lakehouse for IoT Sensor Monitoring}
\textit{Python, Apache Spark, Apache Kafka, PostgreSQL, Parquet, Grafana, Docker}
\begin{itemize}
  \item Engineered a streaming data pipeline that processes 50,000 IoT sensor readings per minute using Kafka producers and consumers with exactly once semantics.
  \item Designed a medallion architecture with Bronze, Silver, and Gold layers stored in Parquet format. Enabled 40\% faster analytical queries compared to flat ingestion.
  \item Implemented schema on read with flexible Delta style partitioning by date and sensor type. Supports ad-hoc queries without pre-defined schemas.
  \item Built Grafana dashboards tracking 20 KPIs including throughput, latency, error rate, and data quality score. Queries run in under a second using materialized Gold layer aggregations.
  \item Added a data quality monitor that detects anomalies such as missing windows, schema drifts, and value outliers using rolling statistics with the IQR method. Auto-pauses ingestion when quality is breached.

% === OPEN SOURCE EXPERIENCE ===
\section{Open Source Experience}
\textbf{Open Source Contributor -- Kestra}\hfill Oct 2025 -- Present\\
\href{https://github.com/kestra-io}{github.com/kestra-io}
\begin{itemize}
  \item Developed the SAP HANA JDBC Subplugin (plugin-jdbc-hana) and enhanced the Google Cloud Platform Dataflow Plugin for Kestra, designing database connectivity and stateful trigger APIs.
  \item Implemented stateful trigger deduplication and custom execution termination hooks (kill/stop lifecycles) for GCP Dataflow tasks, preventing server-side resource leaks and billing overhead.
  \item Currently developing stateful task orchestration features and process termination hooks for additional GCP services (such as Pub/Sub and BigQuery) to minimize runtime overhead.
  \item Contributed 9 merged PRs, authored comprehensive unit and integration test suites, and collaborated asynchronously with distributed maintainers through structured code reviews.
\end{itemize}

% === EDUCATION ===
\section{Education}
\textbf{Master of Computer Applications (MCA)}\hfill Sep 2023 -- Present\\
Kumaraguru College of Technology, Coimbatore \hfill \textbf{CGPA: 7.8 / 10.0}
\\
\textbf{Bachelor of Engineering (B.E) Electronics and Communication Engineering}\hfill Apr 2023 \\
Adithya Institute of Technology \hfill \textbf{CGPA: 7.4 / 10.0}

% === ACHIEVEMENTS ===
\section{Achievements \& Activities}
\begin{itemize}
  \item \textbf{Hacktoberfest 2025 -- Super Contributor:} Multiple accepted PRs across open-source repositories.
  \item \textbf{MCA Department Association Member:} Organized technical events, workshops, and peer-learning sessions.
  \item \textbf{Pear Educator:} Delivered sessions on Git, GitHub, Open Source Contributions, and Data Structures for fellow students.
\end{itemize}

% === ATS NOTE ===
\vfill\noindent\scriptsize
\textbf{ATS Score Estimate:} \textgreater 85/100. Keywords: ETL, Airflow, Kafka, Spark, PostgreSQL, Data Pipeline, Data Warehouse, Lakehouse, Docker, Pandas, SQLAlchemy, Grafana, Parquet, Monitoring, Data Quality. Standard single-column layout, hyperref URLs, clean fonts.

\end{document}
